\begin{center}
{\bf Điếu văn truy điệu Giáo sư Phan Đình Diệu}
\end{center}

Kính thưa quý vị thân hữu, bạn bè, đồng nghiệp,

Thưa gia quyến Thầy Phan Đình Diệu, 

Hôm nay chúng ta cùng nhau ở đây để tiễn đưa Giáo sư Phan Đình Diệu – một nhà khoa học - một nhà giáo - một trí thức lớn của Việt Nam – người đã đặt nền móng cho sự hình thành và phát triển ngành khoa học tính toán và tin học Việt Nam, người có sức ảnh hưởng lớn đến xã hội qua những đóng góp tâm huyết với tầm nhìn xa rộng cho sự phát triển đất nước. 
Sinh năm 1936 tại xã Trung Lộc, huyện Can Lộc, tỉnh Hà Tĩnh trong một gia đình có truyền thống hiếu thuận và hiếu học, ngay từ nhỏ cậu bé Phan Đình Diệu đã bộc lộ một tư chất thông minh đĩnh ngộ. Mới 4 tuổi, cậu đã được thầy giáo cõng đi học cùng các anh chị lớn, và ngay khi 9 tuổi đã được thay mặt thiếu nhi diễn thuyết trước cả hội trường lớn vào những ngày tháng 8 năm 1945 lịch sử. Cha cậu là thầy giáo thường đi dạy xa nhà, mẹ làm nông tần tảo nuôi 7 anh em, là con trưởng nên ngay từ nhỏ, Phan Đình Diệu đã có ý thức tự lập và trách nhiệm, luôn có tấm lòng giúp đỡ và dìu dắt những lớp đàn em đi sau mình. 

Những năm tháng học ở trường cấp III Phan Đình Phùng, nay là trường Trung học phổ thông Phan Đình Phùng, Thành phố Hà Tĩnh, Giáo sư Phan Đình Diệu đã có những người bạn tâm giao, những người sau này trở thành những nhà khoa học trong nhiều lĩnh vực. Ngôi trường ấy cũng cho ông ký ức đẹp về những người thầy mà ông đã thể hiện sự kính yêu qua những bài viết có sức lay động sâu sắc. 

Năm 1954, ông ra Hà Nội và theo học Trường Đại học Sư Phạm Hà Nội. Trong thời gian này, ngoài việc miệt mài học tập, ông cũng đã tự học tiếng Anh và tiếng Pháp, và rất say mê tìm hiểu về Triết học. Chính ở đây ông đã tìm thấy niềm đam mê với Toán học, đã tốt nghiệp thủ khoa xuất sắc. Ngay sau đó, ông được giữ lại trường làm giảng viên và giảng dạy tại đây từ năm 1957 đến 1962.

Năm 1962 ông được cử sang Liên xô làm Nghiên cứu sinh. Năm 1965, với những kết quả xuất sắc khi bảo vệ luận án Phó Tiến sĩ, Ông đã được Hội đồng đề xuất để thực hiện tiếp luận án Tiến sĩ khoa học. 

Năm 1967, ông đã bảo vệ thành công luận án Tiến sĩ khoa học khi ở tuổi 30. Các công trình khoa học có tính đột phá trong giai đoạn này của ông đã được xuất bản trên các tạp chí hàng đầu của Nga, góp phần mở một nhánh mới trong Toán học về Giải tích hàm kiến thiết và được in riêng thành một tập của tạp chí khoa học Steklov của Nga. Năm 1974, được Hiệp hội Toán học Mỹ dịch và xuất bản thành cuốn sách bằng tiếng Anh.

Sau khi về nước, Tiến sĩ khoa học trẻ Phan Đình Diệu làm việc ở Uỷ ban khoa học kỹ thuật nhà nước, tham gia vào Phòng Toán học tính toán và đồng thời kiêm nhiệm tổ trưởng bộ môn Điều khiển học, một trong bốn bộ môn của Viện Toán học trong những ngày đầu thành lập. 

Ngay từ đầu những năm 1970, nhận định về tình hình thế giới, cho rằng thông tin, tri thức sẽ ngày càng đóng vai trò quan trọng và coi đây là cuộc cách mạng lớn thay đổi nhân loại, ông cùng đồng nghiệp đề xuất dự án thành lập Viện Khoa học Tính toán và Điều khiển. 

Năm 1977, Viện Khoa học Tính toán và Điều khiển chính thức được thành lập, và ông là Viện trưởng đầu tiên. Viện là tiền thân của Viện Tin học, sau đổi tên là Viện CNTT hiện nay. 

Giáo sư Phan Đình Diệu đã luôn phát hiện, khuyến khích những cán bộ nghiên cứu trẻ có năng lực, xây dựng các những nhóm nghiên cứu mạnh. Dưới sự lãnh đạo tâm huyết và quyết đoán của ông và với sự quyết tâm và làm việc hăng say các cán bộ của Viện, năm 1979, chiếc máy vi tinh đầu tiên được thiết kế và sản xuất tại Việt Nam và Đông Nam Á. Ông cũng đã lãnh đạo nhóm nghiên cứu hợp tác với Bộ Nội Vụ và Ban Cơ yếu Chính phủ để phát triển các thiết bị tự động giải mật và mật mã hoá. Năm 2007, các nghiên cứu này đã được trao tặng giải thưởng Nhà nước về KH\&CN.

Giáo sư Phan Đình Diệu đã tích cực xây dựng nhiều hợp tác khoa học với quốc tế. Ông đã được mời đến giảng bài hoặc làm Giáo sư mời nhiều trường đại học uy tín trên thế giới ở Pháp, Mỹ, Đức, Nga, Canada, Tây Ban Nha, Nhật Bản, Hungary, Bugary, Ba Lan, Phần Lan. Với ảnh hưởng và uy tín của mình, Ông đã tổ chức thành công các Hội nghị Tin học quốc tế đầu tiên ở Việt Nam, và mời được các chuyên gia Tin học nước ngoài về nước để phát triển các nhóm nghiên cứu trong nước. 

Cùng với sự phát triển của đội ngũ cán bộ KH\&CN trong cả nước, ông đã sáng lập ra Hội Tin học Việt Nam vào cuối năm 1988 và đảm nhận trách nhiệm Chủ tịch Hội hai khoá đầu tiên. 

Vào đầu những năm 1990, ông đã dồn tâm huyết để thực hiện nhiệm vụ được Nhà nước giao biên soạn ``Chính sách phát triển và ứng dụng CNTT'' làm nền tảng cho sự ra đời của Chương trình quốc gia về CNTT. 

Năm 1993, Ông được chỉ định làm Phó trưởng Ban thường trực Ban chỉ đạo Chương trình quốc gia về CNTT và là người trực tiếp soạn thảo Chương trình Quốc gia về CNTT giai đoạn 1995-2000, đặt nền tảng cho việc ứng dụng CNTT vào mọi mặt kinh tế xã hội, như xây dựng chính phủ điện tử, đưa Internet vào Việt Nam, xây dựng các Khoa CNTT trọng điểm và mở các chương trình đào tạo tại các trường đại học.

Dù ở bất cứ cương vị công tác nào, Viện trưởng Viện Tính toán và Điều khiển, Viện phó Viện Khoa học Việt Nam, Chủ tịch Hội đồng chức danh Giáo sư ngành Toán, Trưởng ban xét duyệt và hỗ trợ của Quỹ VIFOTEC của Liên hiệp các Hội KHKT Việt Nam hay Chủ tịch Hội đồng tư vấn về Tin học của Liên hiệp Hội, Giáo sư Phan Đình Diệu luôn thể hiện tinh thần trách nhiệm và sự tâm huyết với công việc.

Ở tuổi lục tuần, Giáo sư Phan Đình Diệu lại trở về giảng dạy tại Trường Đại học Công Nghệ - Đại học Quốc gia Hà Nội.

Trong suốt sự nghiệp dạy học, ông đã dạy nhiều môn học, biên soạn nhiều giáo trình, đặc biệt là các môn học có tính chất khai mở, dự báo những hướng đi mới của sự phát triển khoa học tính toán, như các sách về Logic toán - cơ sở toán học của tin học, Lý thuyết độ phức tạp tính toán, Lý thuyết mật mã và an toàn thông tin. Những bài viết khoa học, những cuốn sách và giáo trình đều được ông viết rất trong sáng, sâu sắc, chặt chẽ và đầy tính sư phạm.

Không chỉ là một nhà khoa học, một giảng viên, suốt cả cuộc đời Giáo sư Phan Đình Diệu còn là một người trí thức đóng góp trí tuệ và tâm huyết cho đất nước. 

Là Đại biểu Quốc hội khoá 5 và khoá 6, tham gia Uỷ viên của Uỷ ban Đối ngoại, rồi Uỷ ban Văn hoá và Giáo dục của Quốc hội, với tư duy độc lập, logic, kết hợp với sự tìm hiểu sâu sắc về triết học và chính trị xã hội, Giáo sư Phan Đình Diệu đã có những cách nhìn sâu rộng về thời cuộc, về xu hướng phát triển và những sự đổi thay cần thiết cho đất nước. Các ý tưởng của ông luôn được hình thành trên cơ sở khoa học, các dữ liệu thực tế và các vấn đề cốt lõi của sự phát triển kinh tế - xã hội; luôn có sự nhất quán, xuyên suốt, mang một tầm nhìn xa, có giá trị qua nhiều thập kỷ. 

Ở mọi diễn đàn trong nước hay quốc tế, ngay cả trên diễn đàn Quốc hội, Mặt trận Tổ Quốc, trước lãnh đạo cấp cao, Phan Đình Diệu luôn thể hiện là một trí thức thẳng thắn, chân thành và tâm huyết vì sự phát triển của đất nước.

Trong xã hội, vốn là người sâu lắng, nhưng tâm hồn chất chứa nét phóng khoáng nghệ sĩ và một tầm hiểu biết xã hội rộng mở, ông có rất nhiều bạn bè tâm giao. Ngôi nhà của vợ chồng ông không chỉ là nơi những bạn bè, đồng nghiệp, học trò trao đổi về học thuật, nơi nhiều trí thức tranh luận về xã hội và đất nước, mà còn là nơi bao bạn bè nghệ sĩ cùng thưởng thức nhiều bản nhạc câu thơ.

Là người con của vùng đất Hà Tĩnh, một nhà khoa học, một giảng viên, một trí thức có nhiều đóng góp với sự phát triển của KH\&CN nước nhà, Giáo sư Phan Đình Diệu luôn trăn trở và giúp đỡ quê hương. Ông là một con người luôn có trách nhiệm và là trụ cột của gia đình, dòng tộc. Là con trưởng, ông trọn lòng hiếu thảo với cha mẹ, chăm lo và dìu dắt sáu người em. Trong gia đình, Giáo sư Phan Đình Diệu là người chồng thuỷ chung, hết mực yêu thương vợ con.

Noi theo tấm gương khoa học trung thực, luôn nhìn thẳng và nói thẳng của đấng sinh thành, cả ba người con ông đều trưởng thành trong sự nghiệp và hiếu thuận với mẹ cha. 

{\em Kính thưa quí vị ! Kính thưa quí quyến !}

Giáo sư Phan Đình Diệu ra đi để lại một khoảng trống cho gia đình, họ tộc, cho quê hương đất nước. Giáo sư đã đi xa nhưng cuộc đời, sự nghiệp, đức độ, thanh danh của Giáo sư mãi còn trong nỗi tiếc thương vô hạn của gia đình, của bạn bè, họ hàng gần xa và của tất cả những ai dù chỉ biết ông qua các trang sách, bài báo. 

Thay mặt Ban tang lễ, xin kính cẩn vĩnh biệt Giáo sư. Xin cầu nguyện cho Giáo sư siêu sinh tính độ. Xin chân thành và thống thiết chia buồn cùng bà quả phụ, các con, các cháu nội ngoại và họ tộc của Giáo sư.

Sau đây, xin tất cả dành một phút mặc niệm Giáo sư Phan Đình Diệu !
(Giáo sư Nguyễn Hữu Đức, phó Giám đốc Đại học Quốc gia Hà Nội đọc. 18-5-2018)
